% !TEX root = ../main.tex
% =========================================================
%  CAPÍTULO 7 — METODOLOGÍA
% =========================================================

\chapter{Capítulo Séptimo - Metodología}
Este capítulo describe el enfoque metodológico adoptado para la gestión y el desarrollo de
TankGo. Se justifica la elección de una metodología ágil iterativa frente a enfoques clásicos
en cascada, se describe el marco de trabajo Kanban implementado mediante Todoist, y se detalla
el uso del control de versiones y las dinámicas de coordinación seguidas a lo largo del
proyecto.

\section{Enfoque de desarrollo: metodología ágil iterativa}

El desarrollo de TankGo se ha abordado siguiendo los principios de las metodologías ágiles
de desarrollo de software~\cite{sommerville2011}, que priorizan la entrega incremental de
valor, la capacidad de respuesta ante cambios en los requisitos y la colaboración continua
con los interesados del proyecto.

Frente a un enfoque en cascada (\textit{waterfall}), en el que las fases de análisis, diseño,
implementación y pruebas se realizan de forma lineal y cerrada, un modelo iterativo resulta
adecuado en proyectos con requisitos emergentes o sujetos a
evolución~\cite{sommerville2011}. TankGo al partir de un prototipo previo y como se ha ido
ampliando el alcance lo largo de las dos fases, este alcance dinámico hace inviable
un planteamiento en cascada y justifica la adopción de un enfoque iterativo e incremental.

\section{Relación con Scrum y elección de Kanban}

Entre las metodologías ágiles más extendidas se encuentran Scrum y Kanban. Scrum estructura
el desarrollo en iteraciones de duración fija denominadas \textit{sprints}, con revisiones
periódicas, \textit{daily stand-up}, revisión de sprint, retrospectiva, y roles diferenciados
como \textit{Product Owner} y \textit{Scrum Master}~\cite{anderson2010kanban}. Este marco está
concebido para equipos multidisciplinares y requiere una coordinación constante entre sus
miembros para que sea efectivo.

TankGo ha sido desarrollado íntegramente por un único estudiante, con supervisión
del director del PFG y colaboración del grupo MORElab. En este contexto de
desarrollo en solitario, las revisiones y artefactos de Scrum resultan desproporcionados y
difíciles de aplicar. Kanban, en cambio, es un marco de gestión
visual del trabajo orientado al flujo continuo de tareas, sin imponer una cadencia
estricta de iteraciones~\cite{anderson2010kanban}. Su principio fundamental es limitar
el trabajo en progreso (\textit{work in progress}, WIP) y hacer visible el estado de
cada tarea en un tablero compartido, lo que favorece la concentración y la detección
temprana de bloqueos. Por estas razones, Kanban se adoptó como marco de referencia
para la gestión del proyecto.

No obstante, el flujo de trabajo adoptado incorpora elementos propios del desarrollo
iterativo que recuerdan a Scrum: para cada módulo o bloque funcional
del sistema se han definido un conjunto de tareas concretas, se han estimado su duración, se
priorizaron y se ejecutaron de forma ordenada hasta su integración y despliegue.
El ciclo por bloque seguía las fases de análisis y diseño,
implementación, pruebas, corrección y despliegue en producción, reproduciendo
el ciclo ágil básico de forma recurrente a lo largo de ambas fases del proyecto.

\section{Gestión de tareas: Todoist como tablero Kanban y agenda diaria}

La herramienta utilizada para la gestión de tareas ha sido Todoist~\cite{Todoist2026},
una aplicación de productividad personal que se ha empleado con una doble función complementaria.

Por un lado, Todoist ha actuado como \textbf{tablero Kanban}: cada módulo funcional del
sistema disponía de su propio bloque de tareas, organizadas por temática
(por ejemplo, \textit{recomendación}, \textit{predicción}, \textit{asistente de voz}).
Dentro de cada bloque se listaban las tareas necesarias (análisis de alternativas,
desarrollo, despliegue en la nube, pruebas, corrección y despliegue definitivo),
a las que se asignaba una prioridad y una estimación de duración. A medida que avanzaba el
trabajo, cada tarea pasaba por los estados \textit{pendiente}, \textit{en progreso}
y \textit{completada}, reproduciendo el flujo de columnas propio de un tablero
Kanban (\textit{backlog} $\rightarrow$ \textit{to do} $\rightarrow$ \textit{in progress} $\rightarrow$ \textit{done}).

Por otro lado, Todoist funcionaba también como \textbf{agenda diaria}: al comienzo de
cada sesión de trabajo se seleccionaban y priorizaban las tareas a abordar ese día,
añadiendo micro-tareas o recordatorios puntales que complementaban la planificación a
nivel de módulo. Esta doble dimensión, tablero de proyecto y planificador diario, ha resultado
útil para compaginar el desarrollo del PFG con otras obligaciones académicas,
permitiendo retomar el trabajo con claridad independientemente del tiempo transcurrido
entre sesiones.

\section{Control de versiones}

El código fuente del proyecto se ha gestionado mediante Git~\cite{git_scm2026,github2026}, con un 
repositorio alojado en GitHub organizado como monorepo. Todo el desarrollo se ha
realizado sobre una única rama principal (\texttt{master}), sin emplear ramas de características
ni flujos de trabajo basados en \textit{pull requests}, decisión coherente con el contexto
de desarrollo en solitario, donde la sobrecarga de un modelo de ramas más elaborado no
aporta beneficio real.

Los \textit{commits} se han realizado de forma frecuente, tras la finalización de
cada tarea o subtarea significativa, lo que permite reconstruir la evolución del
sistema y revertir cambios de forma controlada si fuera necesario. La integración
del repositorio con los pipelines de Cloud Build~\cite{cloud_build2026} cierra el
ciclo de control de versiones: cada \textit{push} a \texttt{master} desencadena
automáticamente la construcción, publicación y despliegue del microservicio correspondiente
en Google Cloud Run, convirtiendo el propio flujo de control de versiones en el disparador
del proceso de entrega continua.

\section{Coordinación y seguimiento}

La supervisión académica y técnica del proyecto se ha organizado en torno a dos canales
de comunicación periódica.

Con el director del PFG se han mantenido reuniones de seguimiento a lo largo de la Fase
2 del proyecto, con una frecuencia aproximada de una reunión cada semana o cada dos semanas, complementadas
con comunicación por correo electrónico para la resolución de dudas puntuales y
la revisión de entregables parciales.

Con el grupo de investigación DEUSTEK/MORElab~\cite{MORElab2026}, a través del proyecto
NextEdge, se ha mantenido una reunión semanal de carácter técnico los viernes, con una
duración aproximada de una hora. Estas sesiones han servido como punto de revisión del
avance del módulo de recomendación, como canal de orientación sobre requisitos y casos
de uso reales, y como motor de algunas de las replanificaciones descritas en el Capítulo 5.
La implicación activa de MORElab ha dotado al proyecto de una dimensión de validación externa
que refuerza su orientación aplicada.